\section{Cyclic Propagation in the All-$\Vdzero$ Cases}
\label{sec:cyclic-all-vd0}

\subsection{The exactly-one-supercritical all-$\Vdzero$ branch}

For a center interval, call the nonempty closed set missed by its two open
endpoint roles a \emph{V-gap}.  It may be a singleton.  If the maximal
closed center trace is $[s,t]$, open coverage gives the weak bounds
$B_i\ge s$ and $A_{i+1}\ge1-t$; thus none of the arguments below discards a
singleton gap.

\begin{proposition}[CE1 one-gap five-map obstruction]
\label{prop:ce1-one-gap}
Assume that every vertex role is $\Vdzero$, $N_+=1$, and the center role is
CE1.  If its boundary trace contains a V-gap, the roles do not cover $H$.
\end{proposition}

\begin{proof}
Use the signed center variables of
Proposition~\ref{prop:signed-center-normal-form}, and put
\[
 w=W,\qquad A=\alpha,\qquad D=\delta,\qquad
 X=R-D,\qquad H=\frac{\eta+A+D}{2R}.
\]
The CE1 sign gives
$D+RA\ge P$ and $RD>wA$, so
\[
 m_3=\min\left\{\frac AR,\frac Dw\right\}=\frac AR.
\]
The exact strict domain needed below is
\begin{equation}
\begin{gathered}
 A>0,\quad D>0,\quad A+wD<P,\quad D+RA\ge P,\\
 A<\frac w2,\quad D<\frac R2,\quad RD-wA>0,
\end{gathered}
\label{eq:ce1-normal-domain}
\end{equation}
Proposition~\ref{prop:signed-one-gap-interface} contains the five actual-V triangle
handoffs, the three duality steps, and the terminal diameter comparison.
It remains only to prove its scalar target
\begin{equation}
 \left(G_{1-D}\circ G_{1-m_3}\circ G_{1-A}\right)(H)>1-X.
 \label{eq:ce1-suffix}
\end{equation}
Here $A<P\le\mu$, because $P=E(1-E)$ and $E\ge\sqrt3/2$.  We first prove
the two input comparisons used at V triangle $4$.  Since $D=R-X$, the inequality
$A+wD<P$ and the identity $Rw=\eta+P$ give
\[
 A<wX-\eta.
\]
Put
\[
 \varphi(q)=wq-\frac{q}{2(1+q)}.
\]
The bound $0<D<R/2$ gives $R/2<X<R$, and $\varphi$ is convex.  At the
endpoints,
\[
 \varphi(R/2)<\frac{Rw}{2}<\eta,
\]
while $\varphi(R)<\eta$ is equivalent to
\[
 E(1-2R^2)<1,
\]
which follows from $E<1$.  Convexity therefore gives
$\varphi(X)<\eta$, and hence
\[
 A<\frac{X}{2(1+X)}.
\]
Using the rational low-root bound in its valid range,
\begin{equation}
 e(A)<\frac{2A}{1-2A}<X.
 \label{eq:ce1-first-low-root}
\end{equation}
We also have
\[
 2R(H-A)=\eta+D+(1-2R)A.
\]
This is positive if $R\le1/2$.  If $R>1/2$, then $A<P=\eta E$ gives
\[
 2R(H-A)>
 \eta\bigl(1-(2R-1)E\bigr)>0.
\]
Thus $H>A$.  The positive signed right trace gives
\[
 \frac{\eta+A+D}{R}<w+D<1,
\]
so $H<1/2<1-A$.

Now consult the exact high-radial catalog for $c_4=1-A$.  The lower and
upper Full ranges would require $H\le A$ and $H\ge1-A$, respectively, so
both are impossible.  On the $L$ or $T_-$ range,
\[
 F_{1-A}(H)\le e(A),
\]
and therefore
\[
 G_{1-A}(H)\ge1-e(A)>1-X
\]
by \eqref{eq:ce1-first-low-root}.  Thus only the
$T_+^{\rm hi}$ range remains.  Its selector gives
\[
 H\le e(A)\le2A+5A^2.
\]
We claim that this range forces $X>1/2$.  Indeed, if $X\le1/2$, the selector
and the lower center inequality give
\[
 2RH=\eta+A+D\ge\eta+P+wA=w(R+A).
\]
Together with $H<2A/(1-2A)$, this implies
\[
 f_R(A)<0,
 \qquad
 f_R(z):=w(R+z)(1-2z)-4Rz.
\]
But $f_R''(z)=-4w<0$ and $f_R(0)=Rw>0$.  We check the other endpoint of
the possible interval for $A$.

If $R\le1/2$, then $0<A<P$.  Direct simplification gives
\[
 f_R(P)
 =(1-E)\left(
 4E^3-2E^2+1-R(2E^3-2E^2+5E)
 \right).
\]
The coefficient of $R$ is positive, so
\[
 f_R(P)\ge
 \frac{1-E}{2}\left(6E^3-2E^2-5E+2\right).
\]
For $\phi(E)=6E^3-2E^2-5E+2$ on
$[\sqrt3/2,1)$,
\[
 \phi'(E)=18E^2-4E-5
 \ge\frac{17}{2}-2\sqrt3>0,
\]
and
\[
 \phi\left(\frac{\sqrt3}{2}\right)
 =\frac{2-\sqrt3}{4}>0.
\]
Hence $f_R(P)>0$.  Strict concavity puts $f_R$ above its endpoint chord,
so $f_R(A)>0$ for $0<A<P$, a contradiction.

If $R>1/2$, the assumption $X\le1/2$ gives
$D\ge R-1/2$, and therefore
\[
 0<A<A_0,
 \qquad
 A_0:=P-w\left(R-\frac12\right)=\frac w2-\eta.
\]
Direct simplification gives
\[
 f_R(A_0)
 =\frac{1-E}{2}(E+11R-3ER-5).
\]
Since $11-3E>0$ and $R>1/2$,
\[
 E+11R-3ER-5
 >\frac{11-3E}{2}+E-5
 =\frac{1-E}{2}>0.
\]
Thus $f_R(A_0)>0$.  Concavity again gives $f_R(A)>0$ on $(0,A_0)$, a
contradiction.  Both ranges are impossible, proving $X>1/2$.

Put $p_1=G_{1-A}(H)$.  If V triangle $3$ is not $T_+^{\rm hi}$, then
$p_1\le A+e(A)<3A+5A^2<29/64<1/2$, and
$2R(H-m_3)=\eta+D-A>\eta-P>0$.  Extensivity gives
$p_1\ge H>m_3$, while $p_1<1/2<1-m_3$; thus both Full ranges are impossible.
On $L$, $T_-$, or the low-$c$ tie one has $F_{1-m_3}(p_1)\le p_1$, so
\[
 G_{1-m_3}(p_1)\ge1-p_1>1-X.
\]
Extensivity of the final reverse map $G_{1-D}$ preserves this strict bound
and proves \eqref{eq:ce1-suffix}.  It remains to treat the selected V triangle-$3$
$T_+$ branch.  Proposition~\ref{prop:reader-universal-tplus} proves once
for all that every selected $T_+$ arc is increasing and strictly concave.

When $d<1-\sqrt3/2$, the arc endpoints are
$(d,d)$ and $(e(d),d+e(d))$, so concavity gives
\begin{equation}
 q\ge p+\frac{d}{e(d)-d}(p-d).
 \label{eq:ce1-selected-chord}
\end{equation}
For $0<d\le\mu$, the valid rational low-root bound gives
\[
 \frac{d}{e(d)-d}>\frac{1-2d}{1+2d}>1-4d>1-5d.
\]
This covers V triangle $4$, because its parameter is $d=A<P\le\mu$.  For the
V triangle-$3$ parameter $d=m_3$, which need not be at most $\mu$, the remaining
high-radial range is handled without the rational estimate.  If
$\mu<d\le1/8$, then
\[
 \frac{d}{e(d)-d}\ge\frac1{1+5d}>1-5d
\]
by $e(d)\le2d+5d^2$.  If
$1/8<d<1-\sqrt3/2$, then $e(d)<(1-d)/2$ and hence
\[
 \frac{d}{e(d)-d}>\frac{2d}{1-3d}>1-5d.
\]
The last comparison is equivalent to $15d^2-10d+1<0$; its smaller root
$(5-\sqrt{10})/15$ is less than $1/8$.  In the low-$c$ V triangle-$3$ range
$1-\sqrt3/2\le d<1/5$, the other endpoint is
$(h(1-d),1-h(1-d))$, and the chord coefficient is
\[
 \frac{1-2h(1-d)}{h(1-d)-d}.
\]
On $3/4\le c\le\sqrt3/2$, direct differentiation of the formula for $h$
shows that $h(c)$ is nonincreasing.  Moreover,
\[
 h\left(\frac34\right)
 =\frac{3}{2(\sqrt6+1)}<\frac7{16}.
\]
Here low $c$ gives $d\ge1-\sqrt3/2>2/15>1/16$, and hence
\[
 h(1-d)<\frac7{16}<\frac{3+d}{7}.
\]
Thus the displayed chord coefficient is greater than $1/3$, while
$1-5d<1/3$.
For $d\ge1/5$, extensivity alone is stronger than the V triangle-$3$ estimate.

Applying these chord bounds first with $d=A$ and then with $d=m_3$ proves,
without an omitted selected branch,
\begin{align*}
 p_1&>L_1:=(2-4A)H-(1-4A)A,\\
 p_2:=G_{1-m_3}(p_1)&>L_2:=(2-5m_3)L_1-(1-5m_3)m_3.
\end{align*}

Finally let
\[
 D_0=P-RA,\qquad
 D_h=A(4R-1)+10RA^2-\eta,\qquad x=Rw.
\]
Feasibility puts $D_0\le D\le D_h$.  We first prove that the actual
feasible value satisfies $D<1/10$.  Since
\[
 P=\frac{xE}{1+E}<\frac x2,\qquad
 \eta=\frac{x}{1+E}>\frac x2,
\]
the contrary assumption $D\ge1/10$ and $A+wD<P$ give
\[
 A<\overline A:=\frac{w(5R-1)}{10}.
\]
The function $D_h(A)$ is increasing, so
\[
 D\le D_h(A)<D_h(\overline A)
 <U(R):=\overline A(4R-1)+10R\overline A^2-\frac x2.
\]
Direct expansion yields
\[
 \frac1{10}-U(R)=-\frac{R}{10}P_4(R),
 \qquad
 P_4(R)=25R^4-60R^3+26R^2+22R-14.
\]
For $y=2R-1\in[0,1]$,
\[
 16P_4(R)=25y^4-20y^3-106y^2+124y-39
 \le-101y^2+124y-39<0.
\]
The last quadratic has discriminant $-380$.  Hence $U(R)<1/10$, a
contradiction, and therefore $D<1/10$.

It remains to compare $L_2$ with $2D+3D^2$.  Put
\[
 J(D)=L_2(D)-\frac{23}{10}D.
\]
Since
\[
 L_2-2D-3D^2-J(D)=3D\left(\frac1{10}-D\right)>0,
\]
it is enough to prove $J(D)>0$.  Exact differentiation gives
\[
 J'(D)=\frac{S(R,A)}{10R^2},\qquad
 S(R,A)=100A^2-(40R+50)A+R(20-23R).
\]
The inequality $D_0\le D_h$ is equivalent to
\[
 x\le A(5R-1+10RA).
\]
Because $A<P<x/2\le1/8$,
\[
 A>L:=\frac{4x}{25R-4}.
\]
Also $\partial_A S=200A-40R-50<-45$, and direct substitution gives
\[
 S(R,L)=-\frac{R q_3(R)}{(25R-4)^2},
\]
where
\[
 q_3(R)=8775R^3-14260R^2+7928R-1120.
\]
For $y=2R-1$,
\[
 8q_3(R)=8775y^3-2195y^2+997y+3007\ge812>0.
\]
Thus $J'(D)<0$ and $J(D)\ge J(D_h)$.

At $D=D_h$ one has $H=2A+5A^2$.  Write $L_{2,h}=L_2(D_h)$.
Expansion gives
\[
 L_{2,h}-A(1+4R)=\frac{A}{R^2}\mathcal Q(R,A),
\]
where
\[
\begin{aligned}
 \mathcal Q(R,A)={}&100A^3R-40A^2R^2-30A^2R
 +12AR^2-15AR+5A\\
 &-4R^3+5R^2-R.
\end{aligned}
\]
On $0\le A\le x/2$,
\[
 \partial_A^2\mathcal Q=20R(30A-4R-3)<0.
\]
Its two endpoint values are
\[
 \mathcal Q(R,0)=x(4R-1)>0,\qquad
 \mathcal Q\left(R,\frac x2\right)=\frac{x}{2}p(R),
\]
where
\[
 p(R)=25R^5-30R^4+20R^3-3R^2-7R+3.
\]
For $y=2R-1$,
\[
 32p(R)=25y^5+65y^4+90y^3+106y^2-35y+5>0,
\]
because the terms of degree at least three are nonnegative and
$106y^2-35y+5$ has discriminant $-895$.  Concavity therefore gives
$L_{2,h}>A(1+4R)$.

On the other hand, $A<P=\eta E$ and $A<P<x/2$ imply
\[
 \frac{D_h}{A}
 =4R-1+10RA-\frac{\eta}{A}
 <C_0:=4R-2+5Rx-\frac x2.
\]
Here we used
\[
 \frac1E>1+\frac x2,
 \qquad
 (1-x)\left(1+\frac x2\right)^2
 =1-\frac{x^2(3+x)}4<1.
\]
Finally,
\[
 1+4R-\frac{23}{10}C_0=\frac{h_3(R)}{20},
 \qquad
 h_3(R)=230R^3-253R^2-81R+112.
\]
Since $h_3''(R)=46(30R-11)\ge184$, strong convexity at $R_0=20/23$
gives
\[
 h_3(R)\ge
 \frac{788}{529}-\frac{(17/23)^2}{368}
 =\frac{289695}{194672}>0.
\]
Consequently
\[
 J(D_h)>A\left(1+4R-\frac{23}{10}C_0\right)>0,
\]
and hence
\[
 p_2>L_2>2D+3D^2.
\]

The low root $e(D)$ is the smaller root of
\[
 H_D(z)=z^2-(1-D)z+(1-D)^2-(1-D)^4,
\]
and direct substitution gives
$H_D(2D+3D^2)=D^2(8D^2+19D-2)<0$.  Hence
$e(D)<2D+3D^2<p_2$.  The one-hit part of
Lemma~\ref{lem:reader-threshold-routing}, applied to the V triangle-$2$ map, gives
\[
 G_{1-D}(p_2)\ge1-e(D)>1-X,
\]
because
\[
 e(D)<2D+3D^2<\frac{23}{100}<\frac12<X.
\]
This is \eqref{eq:ce1-suffix}.  The actual-V triangle and terminal inequalities in
Proposition~\ref{prop:signed-one-gap-interface} now give the contradiction.
\end{proof}

\begin{proposition}[nonempty gap all-$\Vdzero$, $N_+=1$]
\label{prop:nplus-one-all-vd0}
If the center role is CE1 or CE2, every vertex role is $\Vdzero$, and
$N_+=1$, then no cover of $H$ can have a boundary gap in the actual open
vertex-role traces.
\end{proposition}

\begin{proof}
The midpoint argument used above makes $T_0$ uniquely supercritical.  In an
assumed cover, every boundary gap in the vertex-role traces must lie in a
center trace.  CE1 therefore has the one-gap state excluded by
Proposition~\ref{prop:ce1-one-gap}.  For CE2, the exactly-one-gap state is
Proposition~\ref{prop:signed-ce2-one-gap}; the two-gap state is the common
application in Proposition~\ref{prop:reader-common-two-gap}.  This
exhausts the one- and two-gap states, including singleton gaps.  The zero-gap
state is intentionally owned by the center-independent obstruction of
Strategy~4.
\end{proof}
